%------------------------------------------------------------------------------------
%  SECTION 3.1 : FORMULES MATHÉMATIQUES MODERNES AVEC UNICODE-MATH
%------------------------------------------------------------------------------------
\section{Typographie Mathématique Moderne et \textsf{unicode-math}}
\label{sec:formules_mathematiques}

La composition mathématique constitue historiquement le point fort absolu de \LaTeX{}.
Cependant, de nombreuses pratiques héritées des années 1980 persistent encore à tort dans les documents contemporains.
\texttt{ClasseGuillaume} intègre nativement le paquet \textsf{unicode-math} couplé à la police vectorielle mathématique de référence \textbf{XITS Math}, assurant un rendu irréprochable des symboles et une harmonisation d'échelle avec le texte environnant (\texttt{Scale=MatchLowercase}).

\subsection{L'interdiction absolue de \texttt{\$\$...\$\$} et les environnements recommandés}
\label{subsec:interdiction_dollars}

L'une des règles fondamentales de \LaTeXe{} moderne concerne les formules hors-texte :
\begin{boiteSynthese}{Pourquoi la syntaxe \texttt{\$\$...\$\$} est proscrite}
  La syntaxe \texttt{\$\$...\$\$} est une relique primitive du moteur \TeX{} originel (1978).
  En \LaTeX{}, elle perturbe gravement le calcul vertical des marges (\texttt{abovedisplayskip}), court-circuite les options de classes telles que \texttt{fleqn}, et interdit toute numérotation automatique ou référence croisée.
  
  \textbf{Règles à appliquer sans exception :}
  \begin{itemize}
    \item Pour une équation hors-texte non numérotée : utilisez \texttt{\textbackslash [ ... \textbackslash ]}.
    \item Pour une équation hors-texte numérotée : utilisez \texttt{\textbackslash begin\{equation\} ... \textbackslash end\{equation\}}.
    \item Pour des développements multi-lignes alignés : utilisez \texttt{align} ou \texttt{align*}.
  \end{itemize}
\end{boiteSynthese}

\subsection{Ponctuation et texte au sein des équations}
\label{subsec:ponctuation_maths}

Une équation mathématique ne flotte pas dans le vide : elle constitue une phrase grammaticale à part entière ou s'insère dans une proposition existante.
Par conséquent, elle doit impérativement comporter sa \textbf{ponctuation finale} (virgule ou point), placée à l'intérieur de l'environnement mathématique.

De surcroît, les indices ou mots insérés dans une formule ne doivent jamais être saisis en mode mathématique pur, sous peine d'être interprétés comme le produit de variables juxtaposées ($vitesse$ au lieu de $\text{vitesse}$).
On utilise systématiquement la commande \texttt{\textbackslash text\{...\}} du paquet \textsf{mathtools} :

\begin{equation}
  \label{eq:bernoulli_moderne}
  p_1 + \frac{1}{2} \rho v_1^2 + \rho g z_1 = p_2 + \frac{1}{2} \rho v_2^2 + \rho g z_2 + \Delta p_{\text{pertes}},
\end{equation}
où $p$ désigne la pression statique, $\rho$ la masse volumique du fluide, et $\Delta p_{\text{pertes}}$ les pertes de charge singulières et régulières.

\subsection{Démonstration de structures mathématiques complexes}
\label{subsec:structures_complexes}

L'environnement \texttt{align} permet de coordonner des démonstrations pas à pas grâce au point d'alignement \texttt{\&} :
\begin{align}
  \nabla \times \mathbf{E} &= -\frac{\partial \mathbf{B}}{\partial t}, \label{eq:maxwell_faraday} \\
  \nabla \times \mathbf{B} &= \mu_0 \mathbf{J} + \mu_0 \varepsilon_0 \frac{\partial \mathbf{E}}{\partial t}. \label{eq:maxwell_ampere}
\end{align}
Une référence croisée intelligente via \textsf{cleveref} permet d'appeler simultanément les~\cref{eq:maxwell_faraday,eq:maxwell_ampere} sans répétition lourde.

Pour les systèmes différentiels avec conditions aux limites ou fonctions définies par morceaux, l'environnement \texttt{cases} de \textsf{amsmath} garantit un rendu net :
\begin{equation}
  \label{eq:fonction_morceaux}
  f(x, t) = 
  \begin{cases}
    \int_0^t \mathrm{e}^{-\lambda \tau} \cos(\omega \tau) \, \mathrm{d}\tau & \text{si } x \ge 0, \\[2mm]
    \sum_{n=1}^{\infty} \frac{(-1)^n}{n^2 + 1} \sin\left(\frac{n \pi x}{L}\right) & \text{si } x < 0.
  \end{cases}
\end{equation}
Remarquez l'écriture droite du symbole différentiel $\mathrm{d}\tau$ et de la constante d'Euler $\mathrm{e}$, conformément aux normes internationales ISO 80000-2.

Les calculs matriciels s'expriment avec clarté via \texttt{pmatrix} ou \texttt{bmatrix} :
\begin{equation}
  \label{eq:matrice_rotation}
  \mathbf{R}_z(\theta) = 
  \begin{bmatrix}
    \cos\theta & -\sin\theta & 0 \\
    \sin\theta &  \cos\theta & 0 \\
    0          &  0          & 1
  \end{bmatrix},
  \quad
  \det\left(\mathbf{R}_z(\theta)\right) = 1.
\end{equation}
L'\cref{eq:matrice_rotation} illustre la régularité typographique des délimiteurs étirables automatiques.

